\section{Implementation}
\label{sec:implementation}

The production implementation spans four Rust crates
(${\sim}8{,}500$ lines). Rust's zero-cost abstractions let the
trait-based architecture compile away: constraint and backend
polymorphism costs one virtual call per check.

\paragraph{glcore (${\sim}3{,}000$ lines).}
Framework-agnostic core: the \texttt{Constraint} trait
(\texttt{validate(\&self, plan) -> ValidationResult}), the
\texttt{Backend} trait (\texttt{supports\_op}, \texttt{supports\_dtype},
\texttt{supports\_layout}, \texttt{allocate}, \texttt{deallocate},
\texttt{execute\_op}), and the \texttt{Dispatcher} trait; data
structures \texttt{TensorOp}, \texttt{ExecutionPlan},
\texttt{MetricVector} ($d = 5$), \texttt{WeightVector}; and error type
\texttt{GateError} encoding \texttt{ConstraintViolation} (with the
violated constraint and operation), \texttt{NoValidPlan}, and
\texttt{BackendError}.

\paragraph{glproc.}
CPU backend for x86-64 (AVX-512, up to $8\times$ over scalar) and ARM64
(NEON); multi-threaded via \texttt{rayon}; FP32/FP16/INT8; pooled
allocator cutting allocation overhead ${\sim}60\%$ vs.\ per-call
\texttt{malloc}.

\paragraph{glcuda.}
NVIDIA backend on the CUDA runtime API: stream-ordered asynchronous
execution, \texttt{cudaMallocAsync}, FP32/FP16 (Tensor Cores), INT8,
and INT4 (packed kernels); hand-tuned convolution, matmul, attention,
and layer-norm kernels.

\paragraph{glvulkan.}
Cross-vendor GPU backend compiling each operation to SPIR-V compute
shaders dispatched via command buffers; descriptor sets for buffer
binding, pipeline barriers for ordering. Its purpose is portability
across NVIDIA, AMD, Intel, and ARM Mali without CUDA.

\paragraph{Dispatcher and policies.}
The dispatcher routes operations through trait objects
(\texttt{Box<dyn Backend>}), so backends can be registered at runtime;
resource management follows Rust ownership; all fallible paths return
\texttt{Result}. \texttt{ExecutionPolicy} is an enum
(\texttt{Latency}, \texttt{Memory}, \texttt{Energy},
\texttt{Balanced}, \texttt{Custom}) resolved to weight vectors by
constant-time lookup.

\paragraph{Composability.}
Three extension axes: register a new \texttt{Constraint}; register a
new \texttt{Backend}; replace a metric estimator---each without
touching existing components.

\paragraph{Reference implementations for empirical validation.}
Complementing the production crates, two deliberately minimal,
independently coded implementations execute the full
algorithm end-to-end on a real pretrained network and are used in
Section~\ref{subsec:empirical}: \texttt{gate\_resnet}, a
dependency-free Rust binary (im2col convolution, scoped-thread
parallelism, no BLAS or SIMD intrinsics---auditability over speed),
and \texttt{gate\_py}, a NumPy implementation. Both load identical
ImageNet weights, generate real (executable) candidate plans, validate
all seven constraints---including bitwise determinism re-execution---and
select per-policy plans. Their outputs agree with each other and with
three external runtimes to ${\sim}10^{-6}$ relative $L_2$, which is
what qualifies them as the numerical reference for the empirical study.
All code, weights export scripts, and raw results are provided in the
accompanying \texttt{gate-empirical} artifact.
